\documentclass{article}
\usepackage{graphicx} % Required for inserting images

% <added stuff>
\usepackage{amssymb}
\usepackage{amsmath}

\newcommand{\R}{\mathbb{R}}
\newcommand{\N}{\mathbb{N}}
\newcommand{\Z}{\mathbb{Z}}
\newcommand{\Q}{\mathbb{Q}}
\newcommand{\C}{\mathbb{C}}

\newcommand{\power}{\mathcal{P}}
\newcommand{\T}{\mathcal{T}}
\newcommand{\F}{\mathcal{F}}
\newcommand{\G}{\mathcal{G}}
\newcommand{\B}{\mathcal{B}}
\newcommand{\U}{\mathcal{U}}
\newcommand{\mcO}{\mathcal{O}}

\newcommand{\eps}{\varepsilon}
\newcommand{\nullset}{\varnothing}
\newcommand{\lb}{\left(}
\newcommand{\rb}{\right)}
\newcommand{\minus}{\backslash}
\newcommand{\supp}{\text{supp}}
% </added stuff>

\title{Fundamentals 30}
\author{Siddharth K}
\date{June 2025}

\begin{document}

\maketitle

\subsection*{Problem 30.1}
\subsubsection*{Part (1)}
Define
$$f_n(x) = \max(0, 1 - |x - n|).$$
Then this satisfies the constraints.

\subsubsection*{Part (2)}
Let $x \in \R$ and $N > x + 1$. If $n \ge N$, then $f_n(x) = 0$ since $|x - n| = n - x \ge N - x > 1$. So $f_n(x) \to 0$ for all $x$.

\subsubsection*{Part (3)}
Let $\eps < 1$. Then, since $\sup\{d(f(x), g(x)): x \in X\} \ge 1$, $\rho(f, g)$ will never be $< \eps$. So $f_n$ doesn't converge in the uniform topology.

\subsubsection*{Part (4)}
Since $(f_n)$ converges pointwise but not uniformly, the uniform and pointwise topologies aren't equivalent.

\subsubsection*{Part (5)}
Yes, $f_n$ converges in the CC topology. To show this, let $C \subseteq \R$ and $N > \sup(C) + 1$. Then by similar logic as before, if $x \in C$ then $|x - n| > 1$, so $f_n(x) = 0$, and therefore $f_n \to 0$ in the CC topology.

\subsection*{Problem 30.2}
We prove the axioms of a topological basis:
\begin{enumerate}
    \item Let $f \in Y^X$. Choose any $C \subseteq X$ and $\eps > 0$. Then $\rho(f, f) = 0 < \eps$, so $f \in B_C(f, \eps)$.
    \item Let $B_1 = B_{C_1}(f_1, \eps_1)$ and $B_2 = B_{C_2}(f_2, \eps_2)$, and $g \in B_1 \cap B_2$. Let $\delta_1 = \rho(f_1, g) < \eps_1$ and $\delta_2 = \rho(f_2, g) < \eps_2$. Then let $C_3 = C_1 \cup C_2$ (which is compact), $f_3 = g$, and $\eps_3 = \min(\eps_1 - \delta_1, \eps_2 - \delta_2) > 0$. If $h \in B_{C_3}(g, \eps_3)$ then, for each $i = 1, 2$ and $x \in C_i$,
    $$d(f_i(x), h(x)) \le d(f_i(x) + g(x)) + d(g(x), h(x)) < \delta_i + \eps_3 \le \eps_i.$$
    So, for each $i = 1,2$, $h \in B_{C_i}(f_i, \eps_i)$.
\end{enumerate}

\subsection*{Problem 30.3}
First, assume $f_n \to f$ in the CC topology. For every $B_C(f, \eps)$, there is $N$ such that for all $n \ge N$, if $f_n \in B_C(f, \eps)$ then $\rho(f_n, f) < \eps$. This is exactly the same statement as $(f_n)$ converges uniformly to $f$ on $C$. \\

Now, assume that for every $C \subseteq X$, $\rho(f_n, f) \to 0$. Then, given any $B_C(f, \eps)$, by uniform convergence, there is $N$ such that for all $n \ge N$, $\rho(f_n, f) < \eps$, so $f_n \in B_C(f, \eps)$. So $f_n \to f$ in the CC topology.

\subsection*{Problem 30.4}
We have that
$$(f|_C)^{-1}(V) = \{x \in C: f|_C(x) \in V\} = \{x \in C: f(x) \in V\},$$
and
$$f^{-1}(V) \cap C = \{x \in X: f(x) \in V \text{ and }x \in C\} = \{x \in C: f(x) \in V\}.$$
Therefore, the two sets are equal.

\subsection*{Problem 30.7}
Let $x_n \to x$ in $X$ where $(x_n)$ is a sequence in $C$. Since $C$ is closed, $U = X \minus C$ is open. For contradiction, assume $x \notin C$, so that $x \in U$. Since $U$ is open, $U$ is a neighborhood of $x$. But since $x_n \to x$, there exists $N$ such that for all $n \ge N$, $x_n \in U$. This contradicts the fact that all $x_n$ are in $C$, so $x \in C$.

\subsection*{Problem 30.8}
Let $\eps > 0$. Then for any $n$, if $x \ge n\eps$, $f_n(x) \ge \eps$, so it doesn't converge uniformly. Now, let $C \subseteq \R$ be closed and bounded and $\eps > 0$. If we let $N > \frac{\sup(C)}{\eps}$, then if $n \ge N$, then for all $x \in C$,
$$f_n(x) = \frac{x}{n} \le \frac{x}{N} < \frac{x \eps}{\sup(C)} < \eps$$
since $x < \sup(C)$. So $(f_n)$ converges in the CC topology but not uniformly.

\subsection*{Problem 30.9}
The sequence $(f_n)$ converges to $\frac{x}{(1 - x)^2}$ on $(-1, 1)$ in the CC topology but not the uniform topology. \\

Let $C \subseteq (-1, 1)$ be compact. Then there exists $r < 1$ such that $C \subseteq [-r, r]$. For any $x \in C$,
$$|f(x) - f_n(x)| = \left|\sum_{k = n + 1}^\infty kx^k\right| \le \sum_{k = n + 1}^\infty k|x|^k \le \sum_{k = n + 1}^\infty kr^k.$$
Since $|r| < 1$, this sum converges to $0$ as $n \to \infty$. So the sum converges for all $x \in C$. \\

We have that
$$\lim_{n \to 1} f(x) - f_n(x) = \lim_{n \to \infty} \sum_{k = n + 1}^\infty kx^k = \infty,$$
and
$$\sup\{f(x) - f_n(x)\} \ge \lim_{n \to 1} f(x) - f_n(x),$$
so $(f_n)$ is not uniformly convergent.

\end{document}
